\section*{NeurIPS Paper Checklist}

\begin{enumerate}

\item {\bf Claims}
    \item[] Question: Do the main claims made in the abstract and introduction accurately reflect the paper's contributions and scope?
    \item[] Answer: \answerYes{}
    \item[] Justification: The abstract and introduction list four contributions (identification, algorithm, regret bound, experiments) which are stated formally in Theorem~\ref{thm:identification}, Algorithms~\ref{alg:spsc}--\ref{alg:spsc_adaptive}, Theorem~\ref{thm:spsc_regret}, and \S\ref{sec:experiments}; the scope and limitations (e.g.\ approximate-low-rank settings, oracle vs.\ detected change points) are flagged in the relevant sections.

\item {\bf Limitations}
    \item[] Question: Does the paper discuss the limitations of the work performed by the authors?
    \item[] Answer: \answerYes{}
    \item[] Justification: Limitations and extensions are stated in App.~\ref{app:limitations} (scope of the identification theorem, $V_{\rm in}$-dependence of the regret bound, CUSUM detection-delay/false-alarm trade-off, single-rank assumption of Cor.~\ref{cor:rank_adaptive}); the rule-of-thumb crossover $d-r\gtrsim T^{1/6}$ is qualitative, not sharp (\S\ref{sec:exp_phase}).

\item {\bf Theory assumptions and proofs}
    \item[] Question: For each theoretical result, does the paper provide the full set of assumptions and a complete (and correct) proof?
    \item[] Answer: \answerYes{}
    \item[] Justification: All theoretical results (Theorem~\ref{thm:identification}, Theorem~\ref{thm:spsc_regret}, Corollary~\ref{cor:stationary_within}, Corollary~\ref{cor:rank_adaptive}, Proposition~\ref{thm:spsc_adaptive}) state their assumptions explicitly; full proofs are in Appendix~\ref{app:proofs} together with the supporting lemmas (matrix Bernstein, projector concentration, projected nonstationary UCB lemma, prefix concentration).

    \item {\bf Experimental result reproducibility}
    \item[] Question: Does the paper fully disclose all the information needed to reproduce the main experimental results of the paper to the extent that it affects the main claims and/or conclusions of the paper (regardless of whether the code and data are provided or not)?
    \item[] Answer: \answerYes{}
    \item[] Justification: All datasets, dimensions $d$, ranks $r$, segment counts $K$, horizon $T$, seed counts, probe rate, regularisation, window length, and per-baseline hyperparameters are listed in \S\ref{sec:experiments} and Appendix~\ref{app:setup}; per-cell tables in Appendix~\ref{app:percell} report mean$\pm$\,SE for every dataset, $(d,r)$ cell, and method.

\item {\bf Open access to data and code}
    \item[] Question: Does the paper provide open access to the data and code, with sufficient instructions to faithfully reproduce the main experimental results, as described in supplemental material?
    \item[] Answer: \answerYes{}
    \item[] Justification: The supplemental archive contains all experiment scripts (one per figure and per table), per-dataset environment code, cached JSON results, and a README mapping each figure and table to its generator script.

\item {\bf Experimental setting/details}
    \item[] Question: Does the paper specify all the training and test details (e.g., data splits, hyperparameters, how they were chosen, type of optimizer) necessary to understand the results?
    \item[] Answer: \answerYes{}
    \item[] Justification: \S\ref{sec:experiments} states the headline parameters per benchmark; Appendix~\ref{app:setup} gives the full hyperparameter table (probe period, ridge $\lambda$, window $W$, confidence $\delta$, probe cost $c$, baseline-specific knobs) used uniformly across all $(d,r)$ cells without per-dataset retuning.

\item {\bf Experiment statistical significance}
    \item[] Question: Does the paper report error bars suitably and correctly defined or other appropriate information about the statistical significance of the experiments?
    \item[] Answer: \answerYes{}
    \item[] Justification: Every reported number is mean$\pm$\,SE or mean$\pm$\,SEM over at least 10 (typically 10--30) seeds; figures show $\pm 1$\,SE bands; the seed count and aggregation choice are stated in each table caption and in \S\ref{sec:experiments}.

\item {\bf Experiments compute resources}
    \item[] Question: For each experiment, does the paper provide sufficient information on the computer resources (type of compute workers, memory, time of execution) needed to reproduce the experiments?
    \item[] Answer: \answerYes{}
    \item[] Justification: Appendix~\ref{app:setup} reports that the full grid (all benchmarks, 10 seeds, 10 methods) takes 6--20 hours on a single 16-core CPU; no GPU is required (NumPy/SciPy only).

\item {\bf Code of ethics}
    \item[] Question: Does the research conducted in the paper conform, in every respect, with the NeurIPS Code of Ethics \url{https://neurips.cc/public/EthicsGuidelines}?
    \item[] Answer: \answerYes{}
    \item[] Justification: The work is a theoretical and empirical study of a stochastic bandit algorithm on synthetic and publicly available benchmarks; the clinical environments are semi-synthetic surrogates calibrated to published cohort statistics, with no patient-identifiable data, no human subjects, and no clinical deployment.

\item {\bf Broader impacts}
    \item[] Question: Does the paper discuss both potential positive societal impacts and negative societal impacts of the work performed?
    \item[] Answer: \answerNA{}
    \item[] Justification: The paper studies a foundational sequential-decision algorithm with no specific deployment; the positive impact is improved sample efficiency for low-rank decision problems (recommendation, dosing) and negative-impact pathways are not direct. Clinical bandit benchmarks are evaluation surrogates, not deployable dosing policies.

\item {\bf Safeguards}
    \item[] Question: Does the paper describe safeguards that have been put in place for responsible release of data or models that have a high risk for misuse (e.g., pre-trained language models, image generators, or scraped datasets)?
    \item[] Answer: \answerNA{}
    \item[] Justification: The paper releases an algorithm and experiment scripts, not pre-trained models, generative models, or scraped datasets; no high-risk-of-misuse assets are distributed.

\item {\bf Licenses for existing assets}
    \item[] Question: Are the creators or original owners of assets (e.g., code, data, models), used in the paper, properly credited and are the license and terms of use explicitly mentioned and properly respected?
    \item[] Answer: \answerYes{}
    \item[] Justification: All datasets used (UCI Covertype, Pendigits, Satimage; MNIST; Fashion-MNIST; MovieLens-100K; Open Bandit Dataset~\citep{saito2020open}) and all baseline algorithms (LinUCB, D-LinUCB, SW-LinUCB, LowOFUL, VOFUL, LowRank-Reward, LinTS, BOSS, Jedra) are cited at first use; clinical benchmarks follow the IWPC~\citep{international2009estimation} and Rybak et al.~\citep{rybak2020therapeutic} dosing guidelines.

\item {\bf New assets}
    \item[] Question: Are new assets introduced in the paper well documented and is the documentation provided alongside the assets?
    \item[] Answer: \answerYes{}
    \item[] Justification: The supplemental material releases the SPSC and SPSC-Adaptive Python implementations, the per-benchmark environments, all experiment scripts, and a README documenting the figure-to-script and table-to-source mappings.

\item {\bf Crowdsourcing and research with human subjects}
    \item[] Question: For crowdsourcing experiments and research with human subjects, does the paper include the full text of instructions given to participants and screenshots, if applicable, as well as details about compensation (if any)?
    \item[] Answer: \answerNA{}
    \item[] Justification: The paper does not involve crowdsourcing or research with human subjects.

\item {\bf Institutional review board (IRB) approvals or equivalent for research with human subjects}
    \item[] Question: Does the paper describe potential risks incurred by study participants, whether such risks were disclosed to the subjects, and whether Institutional Review Board (IRB) approvals (or an equivalent approval/review based on the requirements of your country or institution) were obtained?
    \item[] Answer: \answerNA{}
    \item[] Justification: The paper does not involve human subjects; no IRB approval is required.

\item {\bf Declaration of LLM usage}
    \item[] Question: Does the paper describe the usage of LLMs if it is an important, original, or non-standard component of the core methods in this research? Note that if the LLM is used only for writing, editing, or formatting purposes and does \emph{not} impact the core methodology, scientific rigor, or originality of the research, declaration is not required.
    \item[] Answer: \answerNA{}
    \item[] Justification: LLMs are not part of the core methodology; the contributions (algorithm, identification result, regret bound, experiments) are independent of any LLM.

\end{enumerate}
